\textbf{结论名称：} 双曲线对称割线的弦与镜像点切线平行

\textbf{一句话：} 从双曲线上同一点 $P$ 出发，画两条关于过 $P$ 的竖直线对称的割线。各自取另一个交点 $M,N$。不管割线怎样成对转动，连线 $MN$ 始终平行于 $P$ 关于 $x$ 轴镜像点 $Q$ 处的切线。

\textbf{条件：} 双曲线 $C:\frac{x^2}{a^2}-\frac{y^2}{b^2}=1$，$a,b>0$；$P(u,v)\in C$。两割线分别为 $y-v=m(x-u)$ 和 $y-v=-m(x-u)$，$m\neq0$；两条直线与 $C$ 除 $P$ 外各有一个\textbf{不同于 $P$ 的有限交点} $M,N$，并且 $P,M,N$ 不共线。直线平行于渐近线时不满足这里的“第二个有限交点”要求。

\textbf{核心结论：} 令 $Q=(u,-v)$，则
\[
\boxed{MN\parallel \text{双曲线在 }Q\text{ 点的切线}}
\]
而 $Q$ 点切线方程为 $\frac{ux}{a^2}+\frac{vy}{b^2}=1$。因此 $MN$ 的方程可写为
\[
\boxed{\frac{ux}{a^2}+\frac{vy}{b^2}=t}\qquad (t\text{ 为某常数}).
\]
若 $v\ne0$，$k_{MN}=-\frac{b^2u}{a^2v}$；若 $v=0$，$MN$ 为竖直线，不讨论它的普通斜率。

\textbf{为什么有用：} 只要认出 $PM,PN$ 关于 $x=u$ 对称，便可不算 $M,N$ 的坐标，直接确定弦 $MN$ 的\textbf{方向}。若还要截距，再补充一个位置条件即可。

\textbf{例：} $C:\frac{x^2}{2}-y^2=1$，$P(2,1)$，$Q(2,-1)$ 的切线 $x+y=1$，所以 $MN$ 的斜率必为 $-1$；但不能直接把 $MN$ 写为 $x+y=1$，因为这只是平行方向，不代表两条线重合。
